\chapter{Book V and the dealers' book: the prints by whose hands they passed through}

\section{Book V, for drawings}

Book V is sixteen sheets and it announces itself: \q{This book for drawings, not water colors.} in lighter pencil on the stationer's thumb-indexed leaf that serves as its index (\lfu{Book V}{index}{16748}). The index is of the book itself: \q{Early Drawings at Rehn Gal. --- left there July 8, 1957. p. 1-5.}; \q{Etchings --- p. 51.}; \q{Loans for travelling shows}; \q{Not yet identified --- but sold --- P. 6}. A leaf holds two to four of Edward Hopper's ink record drawings down the left, a title and a measurement in pencil under each, and beside each one line of hers: buyer, price, the third off, date. There is no exhibition table, no verdict column, no running total.

The drawings went, in June and July 1959 and January 1962, for 150 to 600 dollars: \q{Houses on Skyline 12x18 (“Shacks”)} to George Mladinich, \q{175 - 1/3 = 116 2/3}; Two Big Trees, \q{Cape Ann 1922.}, to Wilson Kierstead at 200; Gloucester House to Dr. B. Nathanson at 600, the highest; \q{Two Nudes - Man + Woman} to Alfredo Valente at 225. The buyers are annotated in pencil afterwards with what she knew of them, \q{friend of John \unc{Lamb}? in big advertising Co. / came here with Grace Root}, \q{Ridgewood, N.J. owns Talbot's House given to Frank Rehn}, and, struck through on the one ruled table, \q{Wm. Ziegler - Indiana Baby Food, reducing pills} with \q{Paid} in the margin (\lf{Book V}{6}{18118}). Two leaves record no transaction at all: they hold the drawings he gave her, inscribed inside the sketch in his hand, \q{To Josie E. H.}, \q{For Jo Edward Hopper}, \q{To Jo from E. Hopper} (\lf{Book V}{7}{17426}; \lf{Book V}{9}{17499}). One is lent: \q{Loaned for Exhibition at National Arts Club - / May 3, 1959. \ldots\ Price \$400. Owned by Jo N. Hopper / for her collection.} The loans page has a single loan, a sanguine, \q{Femme Endormie \ldots\ to Boston Public Library --- Dessins Américains Contemporains / to travel in Europe. Catalogue rec'd for Strasbourg. Sept-Oct. 1956.} and, in pencil, an unclosed \q{(Returned to studio}.

Leaves 51 to 53 are a different document. They are stock-takes of the etchings, and they are the sound of a widow trying to find out where the prints were. \q{Etchings at Rehn Gallery, Oct. 28, 1953.} heads leaf 51 (ref 16364), and under it \q{May 27, 1951. Frank Rehn took from studio in visit there, 11 Etchings :-} followed by nine titles; then a list of what the gallery had before that; then \q{From Kleeman Gallery 1953.}; then, as a heading, a question: \q{Where did \unc{Rehn} gallery get these 4/3 that they sold ?} The transcription warns that leaf 51 is \q{by a long way the densest leaf in the volume} and was read thirteenth in a batch of twelve; a checking pass should start there. Leaf 52 copies a Rehn receipt of 2 July 1959, ten prints at thirty and fifty dollars, \q{340.} less \q{33 1/3 - 113.33} = \q{226.67}, and in the margin: \q{Where did they all come from? Not those owned by Jo \unc{Hopper}, kept in gallery as examples, before E. H. gave up printing any more} (\lf{Book V}{52}{17886}).

The last written leaf records the end of the etchings as a stock: \q{Purchase of Etchings of Edward Hopper by the Philadelphia Museum \unc{of Art.} / Feb. 7'', 1962.} \q{Check signed by Willard P. Graham for \$3500. This in recognition / via Carl Zigrosser of the whole oeuvre of E. H. as etcher - / 162 samples of very early + later work - some in different stages of plate.} and \q{They will have an almost complete collection of the Hopper output of etchings \unc{+ his few drypoints}} (\lf{Book V}{53}{17721}). Book IV writes the same transaction as \q{105 etchings 3500.00}; the counts do not agree, and neither leaf is corrected from the other.

\section{The dealers' book}

The sixth volume is a stationer's cash book, \q{CASH} on the cover, whose front half was used for something else, so that its leaves begin at 51. It is, as its transcription says, \q{indexed by whose hands the work was in}. Leaf 51 closes with two index tables ringed in red pencil: a general one (Downtown Gallery 75 and 82, Whitney Club 78, Babcock 89, Mrs. Sterner 57, Kleemann-Thorman 83, Allison 58, Root and the Utica Museum 59) and one headed \q{Etchings Dealers}: Sterner 88, Weyhe 76, Jennings 80, Keppel 70, Kennedy 85 and 73, Milch 88, Randolph 74, Sidney Phillips 89, Kraushaar 72, Denko 73, and \q{Rehn - See other book}. Each dealer gets a leaf, and the prints move down it: delivered, price, sold, amount received, returned.

Two sentences on the Keppel leaves are the rule for reading every print sale in the archive, and they should be quoted wherever an impression's provenance is built from these books:

\begin{leafquote}
Date of sale = date that payment is rec'd. by the artist. This statement does not refer to the particular copy of the print after which it is entered as the particular / as copies are not numbered + cannot be kept track of. Hence discrepancy in prices - print sold at delivery \hfill (\lf{Dealers}{71}{17690})
\end{leafquote}

\begin{leafquote}
Dates of sale do not belong to the particular print listed - it means any print of that name. \hfill (\lf{Dealers}{72}{17297})
\end{leafquote}

\noindent Hopper did not number his impressions. The ledger row that says a Night Shadows was delivered to Keppel on 12 January 1923 and sold on 5 March 1927 records a print of that name going out and a print of that name being paid for, not the same sheet, and the keeper of the book knew it and wrote it down twice.

The terms are what the volume is for. Weyhe's Book Shop, the earliest dealer entry in the book, took prints on 3 May 1921 and sold Night in the Park and Evening Wind at nine dollars each in November, a 40 per cent discount on fifteen; two of its prints went \q{Gone to Denmark. Mar. 1926.} and one was \q{Gift to Zigrosser} (\lf{Dealers}{76}{16476}). Kennedy took half: \q{Railroad x 2 - Com. 50\% !} and \q{Deux Pigeons - Protested 50\% Com.} (\lf{Dealers}{73}{18070}). The Whitney Studio Club took a tenth. Sterner took 30 per cent. From 1927 the third is general, and on 8 December 1928 the price rose: \q{Increase in prices - 25 to 30 / R. R. Crossing 18 to 25}, with the new figures written over the struck old ones row by row on the Downtown Gallery's leaf (\lf{Dealers}{71}{17690}; \lf{Dealers}{75}{18023}). When a museum bought through a dealer, two commissions came off: \q{Boston Mus. 30 - 15\% Kep. 28 1/2}, and seventeen dollars reached the artist.

The book records, above all, the trouble of keeping track. Edith Halpert's Downtown Gallery returned its consignment in June 1934 and April 1935, and the leaf carries the reconciliation: \q{Account closed, prints returned \ldots\ with the slight discrepancy in inventories - but all prints are accounted for - all have been returned. There has been a substitution of names in its inventories from DownTown Gal. - according to which there would be 2 extra East Side Int. + 1 extra Amer. Landscape. But these replace 1 Lonely House, 1 Night Shadows, 1 Even. Wind, otherwise unaccounted for.} (\lf{Dealers}{82}{17514}). Weyhe's returned its stock while the Hoppers were in Mexico in 1943: \q{Etchings returned to 3 Wash. Sq. in Aug. or Sept. 1943. We in Mexico. package rec'd by Mrs. Waters, Supt. here. \ldots\ Package opened Oct. 17, 43 to find. 4 missing \ldots\ and 1 extra East Side Interior. --- They not supposed to have.} (\lf{Dealers}{77}{16433}). Keppel's, which had become Harlow's, closed in the same year: \q{Rec'd from Harlow, Mar. 25, 43 - / 39 prints. / but on record for Keppel \& Harlow Co. / only 28 prints. / Turn up unaccounted for 11 prints.}; twelve of the thirty-nine went to Allison's, and the leaf ends \q{Account closed.} (\lf{Dealers}{84}{16360}). Kennedy's sold without asking: \q{Lonely House paid only 17, Dec. 8, 53 - we not consulted.}; \q{East Side Interior - June ? 1951. Sold for \unc{\$43} or \$50. We in Mexico, not consulted. \$28 rec'd ?}; and, in the margin, \q{Kennedy got it. Don't know when} (\lf{Dealers}{85}{18068}). The last dealer transaction in the book is Kennedy's, 15 February 1957, and the last date is Zigrosser's visit \q{again ± Mar. 1961. re cataloguing Amer. Etchers / for information very detailed}: the Philadelphia purchase of 1962 was being prepared, from these leaves.

The Rehn leaves at the back, 90 to 95, are \q{the whole back of the book}. They hold the Gloucester and Rockland watercolours of 1926 with one-line descriptions and her verdicts in the date column: \q{Manhattan Bridge entrance - big arch - best of the Lot.} against \q{Painted spring 1926 Brought to R. later in fall - not so good -} (\lf{Dealers}{93}{17434}); \q{Talbot's House (Rockland) - fine white mansard Bgt. by Rehn - Oct. 200 - 1/3 = 133 1/3 Nov. 1926}; and the first oils through the gallery, \q{Fall 1925 750 House by a Railroad - invited by New Society - 600 - 200 Com.} / \q{To Mr. Steven Clark - Mar. 20, 26 paid June 11, 26} (\lf{Dealers}{91}{17125}). On leaf 94, \q{Oils - ready for Hopper}, the Automat stands over a struck \q{The \unc{Cafeteria}}: the picture had another name for a moment.

\section{Her own leaves}

The dealers' book closes, like the etchings collection in Book I, on Josephine Hopper's own work. Leaves 62, 63 and 96 to 98 are her record: the Whitney Studio Galleries' Christmas show of 1929, \q{Total sold \$50. less 25\% = \$37.50 Rec'd Dec. 26, '29}, \q{Sold to Mrs. Force}; and the juries, which is a list of refusals. \q{Submitted to Va. Museum show - Guy du Bois chairman - thrown out.} \q{Artists for Victory at Met. 1941 \ldots\ thrown out.} Corcoran and Chicago in 1943, \q{Invited - E. H. on jury.} \q{Mar. 18'' '46. Turned down by jury of \ill ush Club (thy lowest depths \ill\ Pen \& Brush!} (\lf{Dealers}{97}{16956}). The hole in the leaf at that line is missing paper, the transcriber says, not a hard reading. The last line of the book is hers: \q{Interior of the Ascensions - water color - shown in parish house of the Ascension - late Lent 1952.} (\lf{Dealers}{98}{17161}).

Set beside the tally she made on leaf 56, \q{Sales 34 / Gifts 5 / Dealers 20 / 59}, and the twenty dealers counted by name, these leaves are the difference between the two painters at 3 Washington Square North in a single volume: his prints going out through twenty dealers and coming back with discrepancies, hers going to a jury and coming back thrown out. The Garrulities notebook is the commentary on that difference, and it is the next chapter.
